\documentclass[11pt]{article}
\usepackage[margin=1in]{geometry}
\usepackage{amsmath,amssymb,amsthm,mathtools}
\usepackage{enumitem}
\usepackage{microtype}
\usepackage[hidelinks]{hyperref}

\newtheorem{theorem}{Theorem}[section]
\newtheorem{proposition}[theorem]{Proposition}
\newtheorem{lemma}[theorem]{Lemma}
\newtheorem{corollary}[theorem]{Corollary}
\newtheorem{remark}[theorem]{Remark}
\newcommand{\Pp}{\mathbb P}
\newcommand{\Ee}{\mathbb E}
\newcommand{\1}{\mathbf 1}
\newcommand{\TV}{\mathrm{TV}}
\newcommand{\Poi}{\operatorname{Poisson}}
\newcommand{\Bin}{\operatorname{Bin}}
\newcommand{\Beta}{\operatorname{Beta}}
\newcommand{\LK}{LK}
\newcommand{\fall}[2]{(#1)_{#2}}

\title{The exact connectivity window for bilateral-preference random graphs}
\author{Proof package for the P0456 source question}
\date{September 2026}

\begin{document}
\maketitle

\begin{abstract}
Let independent continuous weights be assigned to the edges of the complete graph $K_n$.
At every vertex keep its $k$ incident edges of smallest weight, and retain an edge in
$\LK(n,k)$ only when both of its endpoints keep it.  Dabirian and Subramanian proved
connectivity and disconnection outside a window of width
$\log\log n\sqrt{\log n}$ around $\log n$ and asked for sharper sufficient conditions
inside that window.  We identify the additive-$O(1)$ transition.  If $Z_n$ is the number
of isolated vertices and $\lambda_n=ne^{-k}$, then, uniformly for every fixed $C<\infty$
and all integers $|k-\log n|\le C$,
\[
 d_{\TV}(\mathcal L(Z_n),\Poi(\lambda_n))\longrightarrow0,
 \qquad
 \Pp(\LK(n,k)\text{ is connected})=e^{-\lambda_n}+o_C(1).
\]
The proof uses an exact marked order-statistic integral for one-root rejection, a canonical
finite interaction/deletion scheme for several roots, a finite Brun--Bonferroni inversion,
and literal elimination of nontrivial components.  No mean-degree heuristic is used.
\end{abstract}

\section{Model and main theorem}

Give every unordered edge $e$ of $K_n$ an independent $U[0,1]$ weight $W_e$.
For a vertex $v$, let $P_v(k)$ be its $k$ incident edges of smallest weight.
The bilateral-preference graph is
\[
 E(\LK(n,k))=\{\{u,v\}:\{u,v\}\in P_u(k)\cap P_v(k)\}.
\]
The continuity of the edge-weight law makes all incident ranks almost surely distinct.
Let $I_v$ be the indicator that $v$ is isolated in $\LK(n,k)$ and put
$Z_n=\sum_v I_v$.

\begin{theorem}[Fixed-window Poisson law and connectivity]\label{thm:main}
For every fixed $C<\infty$, uniformly over integers $k$ satisfying
$|k-\log n|\le C$, with $\lambda_n=ne^{-k}$,
\[
 d_{\TV}(\mathcal L(Z_n),\Poi(\lambda_n))\to0
\]
and
\[
 \Pp(\LK(n,k)\text{ is connected})=e^{-\lambda_n}+o_C(1).
\]
In particular, if $k-\log n\to c$ along an integer-compatible subsequence, then
$Z_n\Rightarrow\Poi(e^{-c})$ and the connectivity probability tends to
$\exp(-e^{-c})$.
\end{theorem}

\begin{corollary}[Sharp source-band threshold]\label{cor:band}
If $k-\log n\to+\infty$, then $\LK(n,k)$ is connected with probability tending to one.
If $k-\log n\to-\infty$, then it is disconnected with probability tending to one.
\end{corollary}

The rest of the paper proves Theorem~\ref{thm:main}.  The key point is that the
rare obstruction at the transition is isolation, whose probability is of order $e^{-k}$;
mean degree alone does not resolve this tail event.

\section{The one-root marked saddle}

We first isolate the analytic input.  For integers $N\ge h$ set
\[
 q_h(x)=\Pp(\Bin(N,x)\ge h),\qquad
 Q_h(x)=\int_0^x q_h(t)\,dt.
\]
If $B_h$ is the $h$th order statistic of $N$ independent uniforms, then
\[
 B_h\sim\Beta(h,N-h+1),\qquad
 q_h(x)=\Pp(B_h\le x),\qquad
 Q_h(x)=\Ee[(x-B_h)_+].
\]
For $m>k$ define
\begin{equation}\label{eq:M}
 M(m,N;k,h)=\frac{\fall m{k+1}}{k!}
 \int_0^1 Q_h(x)^k(1-x)^{m-k-1}\,dx .
\end{equation}
This functional is exactly the probability contribution obtained by exposing the first
$k$ ordered proposals of one root and requiring the selected outside endpoint at deadline
$x$ to have at least $h$ lighter competitors in an independent pool of size $N$.
Indeed, the joint density of the first $k$ order statistics of $m$ uniforms is
$\fall m k(1-x_k)^{m-k}$ on $0<x_1<\cdots<x_k<1$; integrating the first $k-1$
coordinates and integrating by parts gives \eqref{eq:M}.

\begin{lemma}[Uniform clean saddle]\label{lem:saddle}
Fix $C,A,M_0<\infty$.  Suppose
\[
 |k-\log n|\le C,\qquad |m-n|+|N-n|\le Ak,
 \qquad |h-k|\le M_0.
\]
Then, uniformly in these parameters,
\begin{equation}\label{eq:saddle}
 M(m,N;k,h)=e^{-h}\bigl(1+O_{C,A,M_0}(e^{-k/20})\bigr).
\end{equation}
In particular $M(m,N;k,k)=e^{-k}(1+O_C(e^{-k/20}))$.
\end{lemma}

\begin{proof}
Write $\mu=h/(N+1)$ and replace $B_h$ by its mean.  The corresponding transform is
$Q_0(x)=(x-\mu)_+$, and beta integration gives the exact identity
\begin{equation}\label{eq:M0}
 M_0:=\frac{\fall m{k+1}}{k!}
 \int_\mu^1(x-\mu)^k(1-x)^{m-k-1}\,dx=(1-\mu)^m.
\end{equation}
Since $m=n+O(k)$, $N=n+O(k)$ and $h=k+O(1)$,
\[
 M_0=e^{-h}\bigl(1+O(k^2/n)\bigr).
\]
It remains to control the dispersion of $B_h$.

Set $\Delta(x)=Q_h(x)-(x-\mu)_+$.  Convexity of $b\mapsto(x-b)_+$ gives
$\Delta(x)\ge0$.  Interpolate
\[
 Q_\theta(x)=(x-\mu)_++\theta\Delta(x),\qquad0\le\theta\le1,
\]
and let $M_\theta$ be \eqref{eq:M} with $Q_h$ replaced by $Q_\theta$.
Differentiation under the integral yields $M_\theta'\ge0$ and $M_\theta''\ge0$;
hence
\begin{equation}\label{eq:interp}
 0\le M_1-M_0\le M_1'.
\end{equation}
For $x\ge\mu$ write $a=x-\mu$ and
$R(x)=\Ee[(B_h-x)_+]$, so $Q_h(x)=a+R(x)$.
Under the probability measure proportional to
$(x-\mu)^k(1-x)^{m-k-1}\,dx$ on $(\mu,1)$,
\[
 Z:=\frac{x-\mu}{1-\mu}\sim\Beta(k+1,m-k).
\]
Let $\sigma^2=\operatorname{Var}(B_h)$; then
$\sigma=O(\sqrt{k}/n)$ and $\sigma^2/\mu^2\le1/h$.
The normalized version of \eqref{eq:interp} reduces to one lower-boundary term and
one expectation of
$r(x)(1+r(x))^{k-1}$, where $r=R/a$.
Split at
$a_1=\mu k^{-1/3}$ and $a_0=\mu/2$.
For $x\le\mu$ and $0<a\le a_1$, the bounds
$Q_h(x)\le\sigma$ and $R(x)\le\sigma$ give respectively
$\exp[-\tfrac12k\log k+O(k)]$ and
$\exp[-\tfrac13k\log k+O(k)]$.
For $a_1<a\le a_0$, the elementary inequality
$(Y-a)_+\le Y^2/(4a)$ for centered $Y$ gives
$r(x)\le k^{-1/3}$, while the beta/binomial identity and a Chernoff bound give
\[
 \Pp_{\Beta(k+1,m-k)}(a\le a_0)\le e^{-(k+1)/10}.
\]
For $a\ge a_0$, size-biasing the beta density gives
\[
 R(x)\le\mu\Pp(B_h^+>x),\qquad B_h^+\sim\Beta(h+1,N-h+1),
\]
and the order-statistic/binomial identity plus a lower-tail Chernoff bound yields
$r(x)\le2e^{-2(h+1)/35}$.
Substituting these four estimates into \eqref{eq:interp} gives
\[
 0\le M_1/M_0-1=O(e^{-k/20}).
\]
Together with \eqref{eq:M0}, this proves \eqref{eq:saddle}.
\end{proof}

We shall also use the following robust consequence.  It is important that the omitted
ranks may be arbitrary rather than only terminal.

\begin{lemma}[Bounded arbitrary omissions]\label{lem:omit}
Fix $r,C,D<\infty$.  In $r$ one-root marked integrals at $|k-\log n|\le C$, delete at
most $D$ proposal deadlines in each finite collection and change every rejection quota by
at most $D$.  Uniformly over the locations of the deleted ranks, the resulting joint
rejection probability is at most
\begin{equation}\label{eq:omit}
 C_{r,D}k^{O_{r,D}(1)}e^{-rk+C_{r,C,D}k^{2/3}}.
\end{equation}
If the precise total number of deleted deadlines is $d$, the polynomial factor can be
taken to be $C_{r,D}k^d$ after enlarging the $k^{2/3}$ constant.
\end{lemma}

\begin{proof}
Expose the ordered proposal values but not their outside competitor stars.  For a fixed
omission pattern, integrate the retained marked factors in increasing-rank order.
Every missing factor replaces one marked factor by $1$ and contributes at most one choice
among $k$ rank locations.  The remaining terminal integral is bounded by the same
beta/order-statistic saddle with quota shifted by $O(D)$.  Applying the four-sector proof
of Lemma~\ref{lem:saddle} with the shifted quota and summing the finitely many interlacings
of the omitted ranks gives $C_Dk^d$ times an $e^{-k+O(k^{2/3})}$ one-root bound.
Multiplying the $r$ root bounds proves \eqref{eq:omit}.  No favorable-rank assumption is
used.
\end{proof}

\section{Several roots: a canonical finite interaction skeleton}

Fix distinct deterministic roots $R=\{v_1,\ldots,v_r\}$, with $r$ fixed.  Temporarily
delete all root--root edges.  The remaining incident edge families at the roots are
disjoint, so their outside order-statistic arrays are independent.  Write
\[
 X_{a,1}<\cdots<X_{a,k},\qquad U_{a,1},\ldots,U_{a,k}
\]
for their first $k$ outside weights and labels, and $T_a=X_{a,k}$.
Let $S$ be the set of distinct labels among the $rk$ occurrences and
$E=V\setminus(R\cup S)$.

Three kinds of finite defects matter.
\begin{enumerate}[label=(\roman*)]
\item The equality partition of the $rk$ label occurrences has \emph{collision excess}
$c=rk-|S|$.
\item A root--root edge is called endpoint-active if its weight is at most $T_a$ at at
least one endpoint $v_a$; let $p$ be the number of distinct endpoint-active root--root
edges.
\item After the deterministic collision/root preprocessing below, a \emph{support edge}
is a distinct raw edge $\{u,y\}$, with $u\in S$ and
$y\in(R\cup S)\setminus\{u\}$, that is lighter than at least one retained deadline for
label $u$.  Let $q$ be the number of distinct such edges.
\end{enumerate}
Set $K=c+p+q$.

\subsection{Root reinsertion and the two safeguards}

At root $a$, reinsert the deleted root--root edges one at a time.  If currently $q_a(t)$
root edges occupy top-$k$ positions, the selected outside set is exactly
$X_{a,1},\ldots,X_{a,k-q_a(t)}$.  Each successful root-edge insertion expels the current
terminal outside edge.  The $k$th threshold therefore only decreases.  Consequently a
root--root edge heavier than the original $T_a$ can never enter later.
If $d_a$ is the number of endpoint-active root--root edges incident to $a$, the final
root occupancy satisfies $q_a^*\le d_a$, and every outside rank
$1,\ldots,k-d_a$ is certainly retained.  We therefore safely omit the terminal suffix
$k-d_a+1,\ldots,k$.

Two tempting stronger statements are false and are not used.  First, for $r\ge3$, a
root--root edge below both original deleted-array thresholds need not remain in both final
top-$k$ lists after all root edges are reinserted: other root edges may enter first and
lower one endpoint's threshold.  Second, $q_a^*$ need not equal $d_a$.  For example with
$k=2$, outside weights $0.1,0.4$ and two endpoint-active root weights $0.2,0.3$, only one
root edge need occupy the final top two.  The safe suffix rule uses only $q_a^*\le d_a$.

\subsection{Collision canonicalization and support deletion}

Partition proposal occurrences by equal selected label.  In each collision block retain
only the occurrence with the globally smallest deadline, provided it survives the safe
root suffix.  If that global minimum is itself omitted, omit the entire block; never
promote a later occurrence.  Call the resulting retained deadline set $L_0$.

Construct a finite bipartite incidence graph $\Gamma_0$.  Its left vertices are $L_0$.
Its right vertices are the distinct nonexternal support edges just defined, and a right
vertex is adjacent to every retained deadline that it is lighter than.  Let $H$ be the
set of positive-degree right vertices and delete all deadlines in $N(H)$ simultaneously.
The number of right vertices in $H$ is $q$.

\begin{lemma}[Canonical support deletion]\label{lem:canonical}
The construction above is deterministic and order-independent outside the null event of
raw-weight ties.  After the simultaneous deletion, every lighter competitor of every
surviving retained deadline has its other endpoint in $E$.  Consequently the residual
raw-edge stars
\[
 B_u=\{\{u,x\}:x\in E\}
\]
for distinct surviving selected labels $u$ are pairwise disjoint and are disjoint from all
root proposal arrays, root--root edges and charged support edges.  Conditional on the
finite exposed skeleton, these residual stars are independent iid-uniform raw-edge
families.
\end{lemma}

\begin{proof}
The complete incidence graph is determined before any support deadline is deleted.
Deleting left vertices cannot create a new inequality $W_e<D$ and therefore cannot create
new incidences.  Thus simultaneous deletion and any sequential deletion order have the
same surviving left set.  By definition, if a surviving retained deadline had a lighter
competitor with other endpoint in $R\cup S$, that raw edge would be a positive-degree
right vertex and the deadline would have been deleted.  Hence every surviving lighter
competitor goes to $E$.  Distinct selected labels yield distinct unordered edges into $E$;
those edges are also disjoint from the finitely exposed nonexternal family.  Independence
now follows from independence of the original raw weights.
\end{proof}

This is an upper-support device, not an equivalence for joint isolation.  If a root--root
edge enters a root's top $k$, true isolation may impose an additional bilateral-failure
condition.  Dropping such a condition is legitimate only in an upper bound.

\subsection{Defect codimension}

Let $G_L=\{T_a\le L_r k/n\text{ for all }a\}$, with $L_r=8(r+1)$.
For a fixed equality partition with $m=rk-c$ blocks, outside-label permutation symmetry
gives the exact probability
\begin{equation}\label{eq:partition}
 \frac{(n-r)_m}{(n-r)_k^r}=n^{-c+o(1)}.
\end{equation}
There are at most $(rk)^{2c}$ such partitions of collision excess $c$.
On $G_L$, every prescribed endpoint-active root--root edge costs $O_r(k/n)$.
A prescribed selected--selected support edge also costs $O_r(k/n)$.
If a selected label $u$ is absent from a nonowner root's exposed prefix, its edge to that
root has, conditional on the prefix threshold $T$, the exact law $U[T,1]$; therefore its
support activity again costs $O_r(k/n)$.  Distinct charged edges are distinct raw variables.

Split $q=q_{SS}+q_{RT}$ into selected--selected and selected--nonowner-root charges.
The canonical number $d$ of omitted retained deadlines satisfies
\begin{equation}\label{eq:omissions}
 d\le c+2p+2q_{SS}+q_{RT}.
\end{equation}
The selected--selected support graph has at most $2q_{SS}$ incident retained labels; hence
its endpoint choices and omission relief together cost at most $k^{4q_{SS}}$, while its
raw-edge activity contributes $(C_rk/n)^{q_{SS}}$.  The remaining one-unit defect types
cost at most $C_rk^3/n$ after their rank choices and bounded omission relief are included.
Using Lemma~\ref{lem:omit}, for $1\le K\le r$ we obtain
\begin{equation}\label{eq:lowK}
 \Pp(I_{v_1}=\cdots=I_{v_r}=1,\,K,\,G_L)
 \le e^{-rk+C_{C,r}k^{2/3}}
 \left(\frac{C_rk^3}{n}\right)^{c+p+q_{RT}}
 \left(\frac{C_rk^5}{n}\right)^{q_{SS}}.
\end{equation}
For $K\ge r+1$, the same pattern count before using residual rejection gives
\begin{equation}\label{eq:highK}
 \Pp(K\ge r+1,G_L)\le n^{-r-1+o(1)}.
\end{equation}
Finally a binomial Chernoff bound for the $k$th outside order statistic gives
\begin{equation}\label{eq:guard}
 \Pp(G_L^c)\le r\exp[-49(r+1)k/32]=o(e^{-rk}).
\end{equation}
Since $n=e^{k+O_C(1)}$, \eqref{eq:lowK}--\eqref{eq:guard} imply
\begin{equation}\label{eq:positiveK}
 \Pp(I_{v_1}=\cdots=I_{v_r}=1,\,K\ge1)=o_C(e^{-rk}).
\end{equation}

\section{The clean branch and fixed factorial moments}

We next identify the $K=0$ contribution and a matching lower event.
Put $N_0=n-r$.  On global distinctness of all $rk$ outside proposal labels,
$|E|=N_1=n-r-rk$.

\begin{proposition}[Clean upper bound]\label{prop:clean-upper}
Uniformly for fixed $r,C$,
\[
 \Pp(I_{v_1}=\cdots=I_{v_r}=1,K=0)
 \le e^{-rk}\left(1+O_{C,r}(e^{-k/20}+k^2/n)\right).
\]
\end{proposition}

\begin{proof}
On $K=0$, $p=0$ implies that every root's actual top-$k$ list equals its deleted-array
outside list, $c=0$ makes all selected labels distinct, and $q=0$ rules out any lighter
nonexternal support edge.  If a proposal $\{v_a,u\}$ is rejected at $u$, then $u$ has at
least $k$ lighter incident competitors; all of them must therefore lie in $E$.
Conditional on a fixed distinct ordered label array, the stars from selected labels into
$E$ are disjoint and independent.  Hence
\[
 \Pp(I_R=1,K=0)
 \le \frac{(N_0)_{rk}}{(N_0)_k^r}
       M(N_0,N_1;k,k)^r.
\]
The label factor is $1+O_r(k^2/n)$ and Lemma~\ref{lem:saddle} gives the claim.
\end{proof}

For the lower bound let $t_n=8k/N_0$.  Require: (i) all $rk$ outside proposal labels are
distinct; (ii) every root $k$th outside threshold is at most $t_n$; (iii) every root--root
edge has weight greater than $t_n$; and (iv) every selected outside label $u$ has at least
$k$ edges into $E$ lighter than its proposal deadline.  Call this event $L_R$.  It is a
deterministic sufficient event for joint root isolation.

For
\[
 M_t=\Ee\left[\1_{\{X_k\le t\}}\prod_{j=1}^kq_k(X_j)\right],
\]
the same order-statistic integration gives
\begin{align}
 M_t&=\frac{(N_0)_k}{(k-1)!}\int_0^t
 q_k(x)Q_k(x)^{k-1}(1-x)^{N_0-k}\,dx,\label{eq:Mt1}\\
 &=\frac{(N_0)_k}{k!}Q_k(t)^k(1-t)^{N_0-k}
 +\frac{(N_0)_{k+1}}{k!}\int_0^tQ_k(x)^k(1-x)^{N_0-k-1}\,dx.\label{eq:Mt2}
\end{align}
At $t=1$, this is the functional in Lemma~\ref{lem:saddle} with $N=N_1$.
Because all relevant raw-edge blocks are disjoint on distinctness,
\begin{equation}\label{eq:lowerexact}
 \Pp(L_R)=\frac{(N_0)_{rk}}{(N_0)_k^r}
 (1-t_n)^{\binom r2}M_{t_n}^r.
\end{equation}
Moreover
\[
 0\le M_1-M_{t_n}\le\Pp(\Bin(N_0,t_n)\le k-1)\le e^{-49k/16}.
\]
Thus \eqref{eq:lowerexact} and Lemma~\ref{lem:saddle} yield
\begin{equation}\label{eq:lower}
 \Pp(I_{v_1}=\cdots=I_{v_r}=1)
 \ge e^{-rk}\left(1-O_{C,r}(e^{-k/20}+k^2/n)\right).
\end{equation}

Combining \eqref{eq:positiveK}, Proposition~\ref{prop:clean-upper} and
\eqref{eq:lower} proves the central fixed-root estimate
\begin{equation}\label{eq:joint}
 \boxed{\Pp(I_{v_1}=\cdots=I_{v_r}=1)=e^{-rk}(1+o_C(1))}
\end{equation}
uniformly for every fixed $r$ and every deterministic distinct root tuple.

By vertex exchangeability,
\begin{equation}\label{eq:factmom}
 \Ee[(Z_n)_r]=(n)_r\Pp(I_1=\cdots=I_r=1)
 =\lambda_n^r(1+o_C(1))
\end{equation}
for every fixed $r$.

\section{Finite Brun inversion to total variation}

For completeness we record a finite form of the Brun--Bonferroni argument, so no
infinite-moment theorem is hidden in the Poisson step.

\begin{lemma}[Finite Brun bound]\label{lem:brun}
Let $X$ be nonnegative integer-valued, $0\le\lambda\le\Lambda$, and
$\delta_m=|\Ee[(X)_m]-\lambda^m|$, with $\delta_0=0$.  For integers $J\ge0$ and even
$T\ge0$, assuming factorial moments through $J+T+1$ are finite,
\begin{align}
 d_{\TV}(\mathcal L(X),\Poi(\lambda))
 &\le \frac{2\Lambda+\delta_1}{2(J+1)}\\
 &\quad+\frac12\sum_{j=0}^J\frac1{j!}
 \left[\sum_{t=0}^T\frac{\delta_{j+t}}{t!}
 +\frac{\Lambda^{j+T+1}+\delta_{j+T+1}}{(T+1)!}\right].
 \label{eq:brun}
\end{align}
\end{lemma}

\begin{proof}
Set
\[
 B_T^X(j)=\sum_{t=0}^T(-1)^t\frac{\Ee[(X)_{j+t}]}{j!t!}.
\]
The alternating-binomial identity gives, for even $T$,
$B_{T+1}^X(j)\le\Pp(X=j)\le B_T^X(j)$, with bracket width
$\Ee[(X)_{j+T+1}]/[j!(T+1)!]$.  Apply the same brackets to a Poisson variable, whose
factorial moments are $\lambda^m$, compare the brackets, and sum for $0\le j\le J$.
For $j>J$, Markov's inequality bounds the two tails by their first moments.  This gives
\eqref{eq:brun}.
\end{proof}

In the fixed window, $\lambda_n\le e^C$.  Given $\varepsilon>0$, first choose $J$ large
enough to make the first term of \eqref{eq:brun} small, then choose even $T$ large enough
to control the finite Poisson bracket widths, and finally use \eqref{eq:factmom} for the
finitely many required orders.  Hence
\begin{equation}\label{eq:TV}
 d_{\TV}(\mathcal L(Z_n),\Poi(\lambda_n))\to0
\end{equation}
uniformly for $|k-\log n|\le C$.

\section{Literal elimination of nontrivial components}

It remains to prove that isolated vertices are the only asymptotically relevant obstruction
to connectivity.  Let $\tau_v$ be the $k$th incident order statistic at $v$, and set
\[
 a=\frac5{32},\qquad b=\frac{32}{5}.
\]
Let $G$ be the global envelope
\begin{equation}\label{eq:envelope}
 a\frac{k}{n-1}\le\tau_v\le b\frac{k}{n-1}\qquad\text{for every }v.
\end{equation}
The Chernoff rate constants satisfy
\[
 a-1-\log a>1,\qquad b-1-\log b>1.
\]
Since $k=\log n+O(1)$, a union bound over vertices yields
\begin{equation}\label{eq:Gfail}
 \Pp(G^c)=o(1).
\end{equation}

\subsection{Components of size at least seven}

If $S$ is a connected component of size $s$ on $G$, every cut edge from $S$ to $S^c$
has weight at least $ak/(n-1)$: a lighter cut edge would lie below both endpoint kth
thresholds and therefore be retained bilaterally.  Also the retained connected graph on
$S$ contains a spanning tree whose edges all have weight at most $bk/(n-1)$.  Therefore,
for a fixed $S$,
\begin{equation}\label{eq:componentS}
 \Pp(S\text{ is a component},G)
 \le s^{s-2}\left(\frac{bk}{n-1}\right)^{s-1}
 \left(1-\frac{ak}{n-1}\right)^{s(n-s)}.
\end{equation}
For each fixed $7\le s\le12$, summing \eqref{eq:componentS} over $\binom ns$ sets gives
$n^{1-as+o(1)}\operatorname{poly}_s(k)=o(1)$, since $7a=35/32>1$.
For $13\le s\le n/2$, use $n-s\ge n/2$ and $\binom ns\le(en/s)^s$; the resulting
series is bounded by a constant times
\[
 \sum_{s\ge13}n\,[Ck n^{-a/2}]^{s-1},
\]
whose first term tends to zero because $1-13a/2=-1/64<0$ and whose ratio tends to zero.
Thus no component of size between $7$ and $n/2$ survives with high probability on $G$.

\subsection{Components of sizes two through six}

Fix a deterministic set $S$, $2\le s\le6$.  To rank external proposals at $v\in S$,
temporarily delete the $s-1$ internal incident edges.  If
\[
 X_{v,1}<X_{v,2}<\cdots
\]
are the external order statistics and
\[
 H_v(x)=\#\{w\in S\setminus\{v\}:W_{vw}<x\},
\]
then the actual incident rank of the external edge of external rank $j$ is exactly
\begin{equation}\label{eq:rankprofile}
 j+H_v(X_{v,j}).
\end{equation}
Hence the first
\[
 p_0=k-(s-1)
\]
external edges are certainly among the true top $k$, for every realization of the internal
weights.  If $S$ is a component, every one of these forced external proposals must be
rejected at its outside endpoint.

The internal raw sigma-field
$\mathcal F_{\mathrm{int}}(S)=\sigma(W_e:e\subset S)$ is independent of every raw edge
used by the external proposal/rejection/support machinery.  Thus conditioning on the
complete internal rank profile leaves all external laws unchanged.  In particular, if an
outside label $u$ is absent from another component vertex's first $p_0$ external labels
and $T$ is that prefix threshold, then exactly
\[
 W_{bu}\mid\text{ordered prefix}\sim U[T,1].
\]
If this coordinate itself creates a support omission, its $O(k/n)$ activity charge is paid
before revealing and deleting the contaminated deadlines; one never conditions its tail
law on an event that already sees the coordinate.

Applying the same canonical collision/support deletion and Lemma~\ref{lem:omit} to this
fixed-prefix witness, with complexity cutoff $3$, gives uniformly in deterministic $S$
\begin{equation}\label{eq:smallreject}
 \Pp(A_S)\le e^{-sk+o_C(k)}+n^{-3+o(1)},
\end{equation}
where $A_S$ is the event that all forced external proposals are rejected.  The internal
profile removes at most $s(s-1)$ deadlines relative to the full $sk$ proposal array, so
this is precisely a bounded-omission application.

On $G$, a component $S$ has a spanning tree with every internal edge at most $bk/(n-1)$.
For a fixed labelled tree $T$, the light-tree event depends only on internal edges and has
probability at most $(bk/(n-1))^{s-1}$, while $A_S$ depends only on noninternal edges.
They are exactly independent without conditioning on $G$.  Summing over $\binom ns$ sets
and $s^{s-2}$ Cayley trees, \eqref{eq:smallreject} gives, for every $2\le s\le6$,
\[
 \Pp(\exists\ s\text{-component},G)
 \le n^{1-s+o(1)}k^{s-1}+n^{-2+o(1)}k^{s-1}=o(1).
\]
Together with \eqref{eq:Gfail} and the large-component estimate,
\begin{equation}\label{eq:no-components}
 \Pp(\exists\text{ component of size }2,\ldots,\lfloor n/2\rfloor)=o(1).
\end{equation}
The guard failure in \eqref{eq:Gfail} is paid once, after the entire component union.

By \eqref{eq:TV},
\[
 \Pp(Z_n=0)=e^{-\lambda_n}+o_C(1).
\]
A disconnected graph with no isolated vertices has a component of size between $2$ and
$n/2$, so \eqref{eq:no-components} proves Theorem~\ref{thm:main}.

\section{Monotone closure of the source window}

Couple all quotas using the same weight array.  If an edge is retained at quota $k$, it
remains in both endpoints' preferred lists for every larger quota, hence
\[
 \LK(n,k)\subseteq\LK(n,k+1).
\]
Connectivity is therefore increasing in $k$.
If $k_n-\log n\to+\infty$, fix $M>0$ and choose an integer $q_n\le k_n$ eventually with
$M\le q_n-\log n\le M+1$.  Theorem~\ref{thm:main} gives
\[
 \liminf_n\Pp(\LK(n,k_n)\text{ connected})\ge e^{-e^{-M}}.
\]
Letting $M\to\infty$ gives one.  The lower side is analogous: if
$k_n-\log n\to-\infty$, choose $q_n\ge k_n$ with
$-M-1\le q_n-\log n\le-M$; then monotonicity and Theorem~\ref{thm:main} give a limsup at
most $e^{-e^M}$, which tends to zero.  This proves Corollary~\ref{cor:band}.

\section{Scope and source relation}

The theorem answers the connectivity-window question in Section~1.1 of Dabirian and
Subramanian by replacing the source's coarse pairwise no-interaction conditioning with a
canonical finite support-deletion construction.  The source obtains exact conditional
independence after excluding its bad event; here positive interactions are charged once,
every contaminated deadline is deleted, and the remaining external stars are literally
disjoint.  The fixed-window theorem is stronger than the outer-band source bounds and,
by monotonicity, closes the entire source connectivity band.

A bounded literature search through 20 September 2026 did not locate a later primary
publication proving this fixed-width Poisson/connectivity law.  This statement is only a
bounded prior-art record, not a claim of worldwide publication priority.

\begin{thebibliography}{9}
\bibitem{DS}
H. Dabirian and V. Subramanian,
\emph{Connectivity of a family of bilateral preference random graphs},
Journal of Applied Probability \textbf{63} (2026), no.~2, 552--579.
DOI: 10.1017/jpr.2025.10051.
\end{thebibliography}

\end{document}
